\section{Discussion}

The result that most shaped the work is the one in Section \ref{sec:negative}:
tiling alone loses to naive on this card, because the cache hierarchy makes the
naive kernel's redundant traffic cheap. Register blocking, and then feeding the
tensor cores through a swizzled shared layout, are what actually move the number.
Stated as a rule of thumb rather than a measurement, the sequence is that a
technique only pays when the thing it removes is actually costing something, and
the profiler is what tells you whether it is.

The remaining ten percent to cuBLAS at the large shapes, and the much larger gap at
1024 cubed, split cleanly along that line. Table \ref{tab:gap} shows wave
quantization owning the small shape gap, which is a decomposition problem and has
two known answers: split-K and stream-K, both now implemented. The residual at 4096
and 8192 is a different animal, and the hypothesis is per shape kernel selection,
which is what the tile family exists to test. Neither statement is a measurement
yet, and the report says so in the tables rather than in a hedge at the end.

The failed three stage pipeline round is kept in the engineering log as an equal
rank result. A plausible change measured worse and was reverted; the honest record
of that is what makes the changes that did work believable.

One broader observation, offered as an opinion rather than a finding. Most of the
effort in this project went into measurement infrastructure, not kernels: the
locked clock protocol, the allocation gate, the provenance check, the byte stable
report pipeline. That ratio felt wrong while I was in it and looks right in
hindsight, because every claim I had to withdraw in the corrections register was a
claim the infrastructure would have caught before it was published.
